\chapter{Requisitos}

En la siguiente tabla se presentan los requisitos del proyecto. Los requisitos han sido organizados en función de si son requisitos funcionales (RF), de rendimiento (RP), restricciones (RC) y normativos (RN). Aclarar que al no ser aplicado en un sistema real, los requisitos de rendimiento son valores orientativos mas que restricciones que se deben cumplir al 100\%. 

\begin{xltabular}{\textwidth}{l l X}
\caption{Requisitos del proyecto} \label{tab:requisitos} \\
\toprule
\textbf{ID} & \textbf{Categoría} & \textbf{Descripción} \\
\midrule
\endfirsthead

\multicolumn{3}{l}{\small\itshape Continuación de la Tabla~\ref{tab:requisitos}} \\
\toprule
\textbf{ID} & \textbf{Categoría} & \textbf{Descripción} \\
\midrule
\endhead

\endfoot

\bottomrule
\endlastfoot

% --- Funcionales ---
RF01 & Funcional & El sistema debe estimar la actitud (cabeceo y alabeo) y el rumbo de la plataforma en tiempo real. \\
RF02 & Funcional & El sistema debe calcular la altitud respecto al nivel del mar a partir de la presión atmosférica. \\
RF03 & Funcional & El sistema debe estimar la velocidad vertical (tasa de ascenso/descenso) en tiempo real. \\
RF04 & Funcional & Los datos deben transmitirse de forma inalámbrica desde el nodo sensor hasta un receptor. \\
RF05 & Funcional & El receptor debe visualizar los datos en un PFD con disposición Basic-T. \\
\midrule

\pagebreak

% --- Rendimiento ---
RP01 & Rendimiento & Latencia extremo a extremo (sensor $\rightarrow$ PFD) inferior a 200\,ms. \\
RP02 & Rendimiento & Tasa de refresco del PFD $\geq$ 30\,FPS, para una representación fluida. \\
\midrule

% --- Restricciones ---
RC01 & Restricción & Coste total de materiales inferior a 100\,€. \\
RC02 & Restricción & Todo el software utilizado debe ser de código abierto. \\
RC03 & Restricción & La plataforma debe ser compacta y transportable para uso en laboratorio. \\
RC04 & Restricción & Alimentación por USB (laboratorio) o batería. \\
\midrule

% --- Normativos ---
RN01 & Normativa & La disposición y los elementos del PFD deben seguir la normativa aeronáutica aplicable a pantallas electrónicas de vuelo. \\

\end{xltabular}

Los requisitos normativos merecen una aclaración: el sistema no aspira a certificación aeronáutica, pero sí toma las normativas como referencia para que el PFD sea representativo de los instrumentos reales y tenga validez como herramienta educativa.